% GDZ PPN236018647,native600–601;HKUST601–602.
\clearpage\SourceLeaf{600}\ArithmeticHead{571}{R}{1812 FEBR. 29—1814 IUL. 9.}
\NumberedItem{144}
\begingroup\renewcommand{\thefootnote}{[*)]}
\noindent\hspace*{\parindent}Fundamentum theoriae residuorum biquadraticorum \textls[120]{generalis}, per septem\linebreak propemodum annos summa contentione sed semper frustra quaesitum tandem feliciter deteximus eodem die, quo filius\footnotemark\ nobis natus est.
\renewcommand{\thefootnote}{[*)}
\footnotetext{Dieser am 23. Oktober 1813 geborene zweite Sohn aus \textsc{Gauss}' zweiter Ehe mit \textsc{Minna Waldeck} hieß \textsc{Wilhelm}, widmete sich der Landwirtschaft und folgte später seinem älteren Bruder \textsc{Eugen} nach Amerika]}
\par\noindent\hfill 1813 Oct. 23.\quad Gott[ingae].
\par\endgroup

\NumberedItem{145}
\noindent\hspace*{\parindent}Subtilissimum hoc est omnium eorum, quae umquam perfecimus. Vix itaque operae pretium est, his intermiscere mentionem quarumdam simplificationum ad calculum orbitarum parabolicarum pertinentium.
\par\medskip\begingroup\small
\indent Vergl. zu den Nummern 144 und 145 die folgende Stelle aus dem Briefe von \textsc{Gauss} an \textsc{Dirichlet} vom 30. Mai 1828, Werke II, zweiter Abdruck, S. 516: »Die ganze Untersuchung, deren Stoff ich schon seit 23 Jahren vollständig besitze, die Beweise der Haupttheoreme aber (zu welchen das in der ersten Commentation noch \textls[120]{nicht} zu rechnen ist), seit etwa 14 Jahren — (obwohl ich wünsche und hoffe, an letzteren, den Beweisen, noch einiges vereinfachen zu können) — habe ich auf ungefähr drei Abhandlungen berechnet«. Man sehe auch die Bemerkung zu der Nr. 133 sowie den letzten Absatz des Artikels \textbf{22} von \textsc{Bachmanns} Aufsatz, Werke X2, S. 56.
\par\indent In Bezug auf die in der Nr. 145 erwähnten Vereinfachungen zur Berechnung parabolischer Bahnen vergleiche man die Abhandlung \textit{Observationes Cometae secundi a. MDCCCXIII}, Werke VI, S. 25, sowie Werke VII, 1906, S. 338—349.
\par\noindent\hfill\textsc{Klein.\quad Brendel.}
\par\endgroup

\NumberedItem{146}
\noindent\hspace*{\parindent}Observatio per inductionem facta gravissima theoriam residuorum biquadraticorum cum functionibus lemniscaticis elegantissime nectens. Puta si \(a+bi\) est numerus primus, \(a-1+bi\) per \(2+2i\) divisibilis, multitudo omnium solutionum congruentiae
\begingroup\renewcommand{\thefootnote}{[**)]}
\[
1\equiv xx+yy+xxyy\ \text{(mod.\ }a+bi\text{)}\text{\footnotemark},
\]
\renewcommand{\thefootnote}{[**)}
\footnotetext{In der Handschrift steht statt des Kongruenzzeichens \(\equiv\) das Gleichheitszeichen \(=\).]}
\par\endgroup
\noindent inclusis
\[
x=\infty,\quad y=\pm i;\quad x=\pm i,\quad y=\infty,
\]
\noindent fit
\[
=(a-1)^2+bb.
\]
\par\noindent\hfill 1814 Iul. 9.
\par\medskip\noindent\hfill72\textsuperscript{*}
\clearpage\SourceLeaf{601}\ArithmeticHead{572}{L}{TAGEBUCH MIT ERLÄUTERUNGEN.}
\begingroup\small
\indent Die Anzahl Lösungen der Kongruenz \(\text{(mod.\ }a+bi\text{)}\) ist die gleiche wie die der Kongruenz
\[
2\equiv(1+x^2)(1+y^2)\ \text{(mod.\ }p\text{)},
\]
\noindent wo \(p=a^2+b^2\), in reellen ganzen Zahlen (nach \textsc{Dedekind}, Brief an \textsc{Klein}). Man hat also zu suchen, wie groß die Anzahl der Lösungen von
\[
2\equiv u\mathbin{.}v\ \text{(mod.\ }p\text{)}
\]
\noindent ist, bei denen gleichzeitig \(u-1\), \(v-1\) quadratische Reste von \(p\) sind. \textsc{Dedekind} hat für alle Primzahlen \(p<100\) auf diese Weise die \textsc{Gauss}sche Aussage bestätigt gefunden. Andererseits hat R. \textsc{Fricke} (Brief an \textsc{Klein}) darauf hingewiesen, daß die \textls[120]{Gleichung}
\[
1=x^2+y^2+x^2y^2
\]
\noindent die zwischen
\[
x=\sin\operatorname{lemn}u,\quad y=\cos\operatorname{lemn}u
\]
\noindent bestehende Beziehung ist. Der Zusammenhang aber der Theorie der biquadratischen Reste mit den lemniskatischen Funktionen, der durch die Anzahl der Lösungen jener Kongruenz vermittelt wird, bleibt aufzuklären.
\par\noindent\hfill\textsc{Bachmann.}
\par\endgroup
\begin{center}\rule{.18\linewidth}{.3pt}\end{center}

\bigskip
\begin{center}SCHLUSSBEMERKUNG\textsuperscript{*)}.\end{center}
\begingroup\small
\indent Hinter der Nr. 146, mit der die Aufzeichnungen des \textit{Tagebuchs} als solche schließen, sowie auch zwischendurch eingeheftet, finden sich in der Handschrift noch einige Blätter, die mit verschiedenartigen, teils mathematischen, teils nicht mathematischen Aufzeichnungen beschrieben sind\textsuperscript{**)}. Auf der Innenseite der Einbanddecke endlich stehen in eine Falte hineingeschrieben die folgenden Sinnsprüche:
\par\endgroup
\begin{center}
Nil Desperare.\\[.5\baselineskip]
Habeant sibi.\\[.5\baselineskip]
QVA EXEAS HABES.
\end{center}
\noindent\rule{.25\linewidth}{.3pt}
\begingroup\footnotesize
\par\indent *) Diese \textls[120]{Schlußbemerkung} und das folgende \textls[120]{Sachverzeichnis} sind mit einigen geringfügigen Änderungen aus der ersten Ausgabe des \textit{Tagebuchs} übernommen worden.
\par\indent **) Eine dieser Aufzeichnungen mathematischen Inhalts ist oben S. 515 in der Bemerkung zu der Nr. 60 wiedergegeben.
\par\endgroup
\begin{center}\rule{.18\linewidth}{.3pt}\end{center}
